% Generated by scripts/make_paper_tables.py from the frozen results.
% Do not edit by hand; the next build overwrites it.
\begin{table}[tbp]
\centering
\caption{The six notes: underlyings, terms and the figures the issuer discloses}
\label{tab:terms}
\small
\begin{tabular}{l l l r r r r r r}
\toprule
Note & Issuer & Underlyings & Tenor & Coupon & Barrier & Trigger & Est.\ value & Commission \\
\midrule
N1 & JPMorgan & NDX, RUT, SPX & 5.00 & 8.30 & 70.00 & 65.00 & 938.90 & 41.25 \\
N2\textsuperscript{a} & Morgan Stanley & DJX, NDX, RUT & 5.00 & 8.30 & 70.00 & 65.00 & 943.10 & 36.25 \\
N3 & Morgan Stanley & NDX, RUT, SPX & 3.01 & 9.75 & 70.00 & 60.00 & 983.40 & 2.50 \\
N4 & Goldman & NDX, RUT, SPX & 1.92 & 10.15 & 65.00 & 60.00 & 999.00 & 5.25 \\
N5 & BofA & NDX, RUT, SPX & 1.75 & 6.25 & 75.00 & 70.00 & 972.10 & 25.00 \\
N6 & JPMorgan & RUT, SPX, SMH & 1.92 & 12.40 & 65.00 & 60.00 & 967.30 & 22.25 \\
\bottomrule
\end{tabular}
\notetext{\textit{Notes:} Tenor in years from pricing date to stated maturity. Coupon is the contingent rate per annum in per cent, paid only when every underlying is at or above its coupon barrier on an observation date. \textit{Barrier} is that coupon threshold and \textit{Trigger} the principal-repayment threshold, both as a percentage of the initial level. Estimated value and selling commission are per \$1,000 of principal, as disclosed in the pricing supplement. \textsuperscript{a}\,Held out of the ranges quoted in the text: no fitted slice past one year carries a quoted strike at this note's trigger on one leg, so its downside is extrapolated rather than observed. Source: the filings listed in Appendix~\ref{app:terms}.}
\end{table}
